\subsection{Uniform Limits}

% **Role:** convergence of functions and analytic approximation.
%
% Pointwise versus uniform convergence, uniform metric, M-test, preservation of continuity, differentiation, integration.


Uniform convergence asks when limits of functions preserve structure. The key distinction is whether $N$ may depend on $x$, which controls continuity, integration, differentiation, and the supremum metric.

\begin{perspective} \label{per:function-limits-before-function-sequences} \hlt{Function Limits Before Function-Sequence Limits}\\
$\lim_{t\to x}f(t)$ denotes the metric-space function limit of Definition \ref{def:function-limit-metric}.\\
Uniform convergence is the hypothesis that lets the limits in $n$ and $t$ commute.
\end{perspective}

\begin{definition} \label{def:pointwise-convergence-functions} \hlt{Pointwise Limits of Functions}\\
Let $f_n:E\to(Y,d_Y)$.
\begin{enumerate}[label=(\roman*)]
\setlength{\itemsep}{0pt}
\item $f_n\to f$ \hlt{pointwise} on $E$ if $d_Y(f_n(x),f(x))\to0$ for every $x\in E$.
\item If $Y$ is scalar, $\sum f_n\to f$ \hlt{pointwise} on $E$ if $s_N(x)\defeq\sum\limits_{n=1}^{N}f_n(x)\to f(x)$ for every $x\in E$.
\end{enumerate}
No metric on $E$ is required; topology matters only for preserving continuity, limits, or integrals.
\end{definition}

\begin{example} \label{ex:pointwise-limits-do-not-preserve-structure} \hlt{Pointwise Limits Can Lose Structure}
\begin{enumerate}[label=(\roman*)]
\setlength{\itemsep}{0pt}
\item $s_{m,n}\defeq m/(m+n)$ gives $\lim_n\lim_m s_{m,n}=1$ but $\lim_m\lim_n s_{m,n}=0$.
\item For $f_n(x)=x^2/(1+x^2)^n$, $\sum_{n=0}^{\infty}f_n$ equals $0$ at $0$ and $1+x^2$ otherwise, by Theorem \ref{thm:geometric-series-nonnegative}.
\item $f_n(x)=\sin(nx)/\sqrt n$ satisfies $f_n\to0$ pointwise but $f_n'(0)=\sqrt n\to+\infty$.
\item $f_n(x)=n^2x(1-x^2)^n$ on $[0,1]$ satisfies $f_n\to0$ pointwise, while $\int_0^1 f_n\,dx=n^2/(2n+2)\to+\infty$.
\end{enumerate}
\hlt{Key Insight}: Pointwise control can miss global spikes and oscillations.
\end{example}

\begin{figure}[H]
  \centering
  \begin{tikzpicture}[
  >=Stealth,
  every node/.style={font=\small},
  axis/.style={->, gray!70, line width=0.45pt},
  curve/.style={line width=0.75pt},
  limit/.style={gray!70, dashed, line width=0.55pt},
  point/.style={circle, fill=black, inner sep=1.35pt}
]
  \draw[axis] (-0.15,0) -- (4.7,0) node[right] {$x$};
  \draw[axis] (0,-0.15) -- (0,3.0) node[above] {$y$};

  \draw[gray!65] (4,0.05) -- (4,-0.05) node[below] {$1$};
  \draw[gray!65] (0.05,2.4) -- (-0.05,2.4) node[left] {$1$};

  \draw[curve, blue!70!black, domain=0:4, variable=\xx, samples=80] plot (\xx,{2.4*(\xx/4)});
  \draw[curve, dartmouthgreen!75!black, domain=0:4, variable=\xx, samples=100] plot (\xx,{2.4*(\xx/4)^2});
  \draw[curve, jazzberryjam!80!black, domain=0:4, variable=\xx, samples=100] plot (\xx,{2.4*(\xx/4)^4});
  \draw[curve, black!70, domain=0:4, variable=\xx, samples=120] plot (\xx,{2.4*(\xx/4)^8});

  \draw[limit] (0,0) -- (4,0);
  \node[point] at (4,2.4) {};
  \node[right] at (4.1,2.4) {$g(1)=1$};
  \node[below, black!70] at (2.15,-0.02) {$g(x)=0$ for $x<1$};

  \node[blue!70!black] at (2.1,1.55) {$x$};
  \node[dartmouthgreen!75!black] at (2.55,0.82) {$x^2$};
  \node[jazzberryjam!80!black] at (3.08,0.55) {$x^4$};
  \node[black!70] at (3.5,0.32) {$x^8$};
\end{tikzpicture}

  \caption{$x^n\to0$ pointwise on $[0,1)$, but the endpoint prevents uniform convergence on $[0,1]$.}
  \label{fig:xn-pointwise-not-uniform}
\end{figure}

\begin{example} \label{ex:xn-pointwise-not-uniform} \hlt{$x^n$ Converges Pointwise But Not Uniformly}\\
On $[0,1]$, let $g_n(x)=x^n$. Then $g_n\to g$ pointwise, where
$g(x)=
\begin{cases}
0, & 0\leq x<1,\\
1, & x=1.
\end{cases}$.\\
The convergence is not uniform: for each $n$, some $x<1$ has $\abs{x^n-g(x)}=x^n>1/2$.
\end{example}

\begin{figure}[H]
  \centering
  \begin{tikzpicture}[
  >=Stealth,
  every node/.style={font=\small},
  axis/.style={->, gray!70, line width=0.45pt},
  strip/.style={gray!65, dotted, line width=0.65pt},
  limit/.style={black, line width=0.9pt},
  tail/.style={blue!70!black, line width=0.75pt},
  domainbar/.style={black!75, line width=1.3pt}
]
  \draw[axis] (-0.35,0) -- (7.2,0) node[right] {$x$};
  \draw[axis] (0,-0.5) -- (0,3.3) node[above] {$y$};

  \draw[domainbar] (0.75,-0.18) -- (6.45,-0.18);
  \node[below] at (3.6,-0.18) {$E$};

  \draw[strip, domain=0.75:6.45, variable=\xx, samples=120]
    plot (\xx,{2.0 + 0.24*sin(85*\xx) + 0.52});
  \draw[limit, domain=0.75:6.45, variable=\xx, samples=120]
    plot (\xx,{2.0 + 0.24*sin(85*\xx)});
  \draw[strip, domain=0.75:6.45, variable=\xx, samples=120]
    plot (\xx,{2.0 + 0.24*sin(85*\xx) - 0.52});

  \draw[tail, domain=0.75:6.45, variable=\xx, samples=120]
    plot (\xx,{2.0 + 0.24*sin(85*\xx) + 0.18*sin(190*\xx)});
  \draw[tail, blue!45!black, domain=0.75:6.45, variable=\xx, samples=120]
    plot (\xx,{2.0 + 0.24*sin(85*\xx) - 0.15*cos(170*\xx)});

  \node[right] at (6.52,2.78) {$f+\varepsilon$};
  \node[right] at (6.52,2.18) {$f$};
  \node[right] at (6.52,1.68) {$f-\varepsilon$};
  \node[blue!65!black, anchor=west] at (1.25,2.95) {$f_n,\ n\geq N$};
\end{tikzpicture}

  \caption{One index $N$ puts all later graphs inside the same $\varepsilon$-strip around $f$ on $E$.}
  \label{fig:uniform-convergence-epsilon-strip}
\end{figure}

\begin{definition} \label{def:uniform-convergence-functions} \hlt{Uniform Convergence}\\
For $f_n,f:E\to(Y,d_Y)$, $f_n\rightrightarrows f$ on $E$ if $\forall \varepsilon>0$ $\exists$ $N \in \N$ such that $n\geq N,\ x\in E \Longrightarrow d_Y(f_n(x),f(x))\leq\varepsilon$.\\
In the scalar-valued case, take $d_Y(u,v)=\abs{u-v}$; $\sum f_n$ converges uniformly on $E$ if its partial sums do.
\end{definition}

\begin{perspective} \label{per:uniform-convergence-epsilon-strip} \hlt{One Tail, All Points}\\
Pointwise convergence lets $N$ depend on $x$; uniform convergence uses one $N$ for all $x$.\\
Equivalently, late graphs lie in one $\varepsilon$-strip around $f$.
\end{perspective}

\begin{theorem} \label{thm:uniform-cauchy-criterion} \hlt{Cauchy Criterion for Uniform Convergence}\\
For complex-valued $f_n:E\to\C$, $f_n$ converges uniformly on $E$ if and only if
\begin{align}
\forall\varepsilon>0\ \exists N\in\N:\quad m,n\geq N,\ x\in E \Longrightarrow \abs{f_n(x)-f_m(x)}\leq\varepsilon \nonumber
\end{align}
\begin{proof}
\begin{enumerate}[label=\roman*.]
\setlength{\itemsep}{0pt}
\item[$(\Rightarrow)$] If $f_n\rightrightarrows f$, choose $N$ so that $\abs{f_k(x)-f(x)}\leq\varepsilon/2$ for $k\geq N$ and $x\in E$. Then $\abs{f_n(x)-f_m(x)}\leq\varepsilon$.
\item[$(\Leftarrow)$] For each $x$, $(f_n(x))$ is Cauchy in $\C$, so define $f(x)\defeq\lim f_n(x)$.\\
Fix $n\geq N$ and let $m\to\infty$ in the criterion; then $\abs{f_n(x)-f(x)}\leq\varepsilon$ for all $x\in E$, hence $f_n\rightrightarrows f$.
\end{enumerate}
\hlt{Structural Commentary}: The same $N$ gives Cauchy control simultaneously for every $x\in E$.\\
\hlt{Key Insight}: Uniform Cauchy tails are exactly uniform convergence.
\end{proof}
\end{theorem}

\begin{theorem} \label{thm:uniform-convergence-supremum-test} \hlt{Supremum Test for Uniform Convergence}\\
Suppose $f_n(x)\to f(x) \ \forall x\in E$, and $M_n\defeq\sup_{x\in E}\abs{f_n(x)-f(x)}$. Then $f_n\rightrightarrows f$ on $E$ if and only if $M_n\to0$.
\begin{proof}
\begin{enumerate}[label=\roman*.]
\setlength{\itemsep}{0pt}
\item[$(\Rightarrow)$] If $f_n\rightrightarrows f$, then $\forall \varepsilon>0$ $\exists$ $N$ such that $n\geq N$ implies $\abs{f_n(x)-f(x)}\leq\varepsilon$ for all $x\in E$. Hence $M_n\leq\varepsilon$.
\item[$(\Leftarrow)$] If $M_n\to0$, then $\forall  \varepsilon>0$ $\exists$ $N$ such that $n\geq N$ implies $M_n\leq\varepsilon$. Hence $\abs{f_n(x)-f(x)}\leq\varepsilon$ for all $x\in E$.
\end{enumerate}
\hlt{Structural Commentary}: Uniform convergence is convergence in the supremum error.\\
\hlt{Key Insight}: The pointwise errors must vanish as one global quantity.
\end{proof}
\end{theorem}

\begin{theorem} \label{thm:weierstrass-m-test} \hlt{Weierstrass $M$-Test}\\
Let $f_n$ be functions on $E$ and suppose $\abs{f_n(x)}\leq M_n$ for every $x\in E$.\\
If $\sum M_n$ converges, then $\sum f_n$ converges uniformly on $E$.
\begin{proof}
For $m\geq n$ and $x\in E$, $\abs{\sum\limits_{k=n}^{m}f_k(x)}\leq\sum\limits_{k=n}^{m}M_k$.\\
Since $\sum M_k$ converges, Theorem \ref{thm:series-cauchy-criterion} makes the right-hand side $<\varepsilon$ for all large $m,n$.\\
Thus the partial sums of $\sum f_k$ are uniformly Cauchy, so Theorem \ref{thm:uniform-cauchy-criterion} applies.\\
\hlt{Structural Commentary}: Absolute numerical tails dominate functional tails.\\
\hlt{Key Insight}: A summable uniform majorant gives uniform convergence of a function series.
\end{proof}
\end{theorem}

\begin{theorem} \label{thm:uniform-convergence-interchanges-function-limits} \hlt{Uniform Convergence Interchanges Limits}\\
Let $E$ lie in a metric space, let $x$ be a limit point of $E$, and suppose $f_n\rightrightarrows f$ on $E$.\\
If $A_n\defeq \lim\limits_{t\to x}f_n(t)$ exists for every $n$, then $(A_n)$ converges and $\lim\limits_{t\to x} f(t)=\lim\limits_{n\to\infty}A_n$.
\begin{proof}
By Theorem \ref{thm:uniform-cauchy-criterion}, $\forall \varepsilon>0$ $\exists$ $N$ such that $m,n\geq N$ and $t\in E$ imply $\abs{f_n(t)-f_m(t)}\leq\varepsilon$.\\
Letting $t\to x$ gives $\abs{A_n-A_m}\leq\varepsilon$, hence $A_n\to A$ for some $A\in\C$.\\
For any $n$, $\abs{f(t)-A}\leq \abs{f(t)-f_n(t)}+\abs{f_n(t)-A_n}+\abs{A_n-A}$. Given $\varepsilon>0$, choose $n$ so $\abs{f(t)-f_n(t)}<\varepsilon/3$ for all $t\in E$ and $\abs{A_n-A}<\varepsilon/3$; then take $t$ sufficiently close to $x$ so $\abs{f_n(t)-A_n}<\varepsilon/3$. Thus $\lim_{t\to x}f(t)=A$.\\
\hlt{Structural Commentary}: Uniform convergence controls the outer two terms in the three-term estimate.\\
\hlt{Key Insight}: Uniformity is precisely what allows the $n\to\infty$ limit to move past the $t\to x$ limit.
\end{proof}
\end{theorem}

\begin{theorem} \label{thm:uniform-convergence-preserves-continuity}\label{thm:uniform-convergence-main-source-todo} \hlt{Uniform Limits Preserve Continuity}\\
If $(f_n)$ is a sequence of continuous functions on $E$ and $f_n\rightrightarrows f$ on $E$, then $f$ is continuous on $E$.
\begin{proof}
Let $x\in E$. If $x$ is isolated, $f$ is continuous at $x$.\\
If $x$ is a limit point, then $\lim\limits_{t\to x}f_n(t)=f_n(x)$ for all $n$, so Theorem \ref{thm:uniform-convergence-interchanges-function-limits} gives $\lim\limits_{t\to x}f(t)=\lim\limits_{n\to\infty}f_n(x)=f(x)$.\\
\hlt{Structural Commentary}: Uniform convergence passes continuity to the limit.\\
\hlt{Key Insight}: Uniform convergence protects continuity from the failure modes in Example \ref{ex:pointwise-limits-do-not-preserve-structure}.
\end{proof}
\end{theorem}

\begin{theorem} \label{thm:dini-monotone-uniform-convergence} \hlt{Dini's Monotone Uniform Convergence Theorem}\\
Let $K$ be compact and let $f_n,f:K\to\R$ be continuous. If $f_n(x)\downarrow f(x)$ for every $x\in K$, then $f_n\rightrightarrows f$ on $K$.
\begin{proof}
Set $g_n\defeq f_n-f$. Then $g_n$ is continuous, $g_n\downarrow0$, and $g_n\geq0$. Given $\varepsilon>0$, let $K_n\defeq\{x\in K:g_n(x)\geq\varepsilon\}$. Then $(K_n)$ is nested compact and $\bigcap_n K_n=\emptyset$, so Corollary \ref{cor:nested-nonempty-compact-intersection} gives $K_N=\emptyset$ for some $N$.\\
Thus $0\leq f_n(x)-f(x)<\varepsilon$ for all $x\in K$ and $n\geq N$.\\
\hlt{Structural Commentary}: Monotone convergence turns the bad sets into a nested compact sequence.\\
\hlt{Key Insight}: Compactness turns pointwise eventual escape from every bad set into a uniform escape.
\end{proof}
\end{theorem}

\begin{definition} \label{def:supremum-norm-continuous-functions} \hlt{The Supremum Metric on Bounded Continuous Functions}\\
For a metric space $X$, let $C_b(X)$ be the bounded complex-valued continuous functions and\\
define $\norm{f}_{\infty}\defeq\sup_{x\in X}\abs{f(x)}$; $d(f,g)\defeq\norm{f-g}_{\infty}$.\\
Then $d$ is a metric on $C_b(X)$, induced by the norm of Definition \ref{def:norm-and-normed-vector-space}; $d$-convergence is uniform convergence by Theorem \ref{thm:uniform-convergence-supremum-test}. If $X$ is compact, then $C(X)=C_b(X)$ by Corollary \ref{cor:continuous-rk-image-compact-closed-bounded}.
\end{definition}

\begin{perspective} \label{per:supremum-metric-function-space-lens} \hlt{A Metric Space Whose Points Are Functions}\\
Uniform convergence is sequence convergence in the metric space $C_b(X)$.
\end{perspective}

\begin{perspective} \label{per:four-limit-notions-function-space-view} \hlt{Four Limit Notions}\\
The four limit notions are sequence limits, function limits, pointwise function limits, and uniform function limits. By Theorem \ref{thm:uniform-convergence-supremum-test}, the last is ordinary sequence convergence in the supremum metric.
\end{perspective}

